% Q3: T-equivalent circuit for coupled inductors (La, Lb series; Lc shunt)
\documentclass[border=10pt]{standalone}
\usepackage[american voltages, american currents]{circuitikz}
\usepackage{amsmath}
\renewcommand{\familydefault}{\sfdefault}

\begin{document}
\begin{circuitikz}[
    line width=0.8pt,
    every node/.style={font=\sffamily}
]

% Left port terminal (top)
\draw (0,3) -- (1,3);
\fill (0,3) circle (2pt);
\node[left] at (0,3) {1};

% La (series inductor on port 1 side)
\draw (1,3) to[L, l=$L_a$] (3.5,3);

% Junction node
\fill (3.5,3) circle (2pt);

% Lb (series inductor on port 2 side)
\draw (3.5,3) to[L, l=$L_b$] (6,3);

% Right port terminal (top)
\draw (6,3) -- (7,3);
\fill (7,3) circle (2pt);
\node[right] at (7,3) {2};

% Lc (shunt inductor from junction to bottom rail)
\draw (3.5,3) to[L, l_=$L_c$] (3.5,0);

% Bottom rail
\draw (0,0) -- (7,0);
\fill (0,0) circle (2pt);
\fill (7,0) circle (2pt);
\node[left] at (0,0) {1'};
\node[right] at (7,0) {2'};

% Port labels
\node[above left] at (0.5,3.2) {Port 1};
\node[above right] at (6.5,3.2) {Port 2};

% Current arrows
\draw[->, >=stealth, thick] (0.3,3.4) -- (1.5,3.4)
    node[midway, above] {$i_1$};
\draw[->, >=stealth, thick] (6.7,3.4) -- (5.5,3.4)
    node[midway, above] {$i_2$};

\end{circuitikz}
\end{document}
